\documentclass[10pt,a4paper]{amsart}
\usepackage[margin=19mm]{geometry}
\usepackage{amsmath,amssymb,mathtools}
\usepackage[scr=rsfs,cal=euler]{mathalfa}
\usepackage{xfrac}
\usepackage[unicode,hidelinks,bookmarks=false]{hyperref}
\newcommand{\ie}{i.e.\ }
\newcommand{\cf}{cf.\ }
\newcommand{\resp}{respectively\ }
\newcommand{\Subsection}{\subsection}
\providecommand{\nobreakdash}{\nobreak\hbox{-}}
\setlength{\emergencystretch}{2em}
\multlinegap0pt
\usepackage[shortlabels]{enumitem}
\usepackage{xcolor}
\usepackage{amssymb}
\usepackage{mathtools}
\usepackage{dsfont}
\usepackage{float}
\usepackage{algorithm}\usepackage{algpseudocode}
\newtheorem{assumption}{Assumption}
\newtheorem{lemma}{Lemma}
\newcommand{\R}{\mathbb{R}}      %Numeros reales
\title{On Runge-Kutta convolution quadrature based fractional variational integrators}
\author{Khaled Hariz Belgacem, Sina Ober-Bl\"obaum, Fernando Jim\'enez}
\date{Source: arXiv:2606.23302v1 (CC BY 4.0)}
\begin{document}
\maketitle

\noindent\emph{Source and licence.} On Runge-Kutta convolution quadrature based fractional variational integrators. Version arXiv:2606.23302v1 (CC BY 4.0), distributed by arXiv under the Creative Commons Attribution 4.0 International licence, http://creativecommons.org/licenses/by/4.0/, as recorded in the arXiv licence metadata for this exact version and shown on the abstract page https://arxiv.org/abs/2606.23302v1. Full licence notice: \texttt{licenses/arxiv-2606.23302-CC-BY-4.0.txt}.\par\smallskip

\noindent\textit{Editorial scope.} Counted unit: the article's lemma lemma:Semi\_group (source0.tex lines 492--500). The article's complete original proof of it is delivered with it (source0.tex lines 501--568, one \texttt{proof} environment of 68 lines). No mathematics is written, repaired or re-derived: the statement and the proof are the article's own lines, in the article's own notation, and no retained line is substituted or re-flowed.  Retained uncounted as the source-local context the proof needs: lines 296--303.  Nothing was omitted from any retained span, so the package carries no omission record.  No retained line contains a citation, so the wrapper carries no bibliography and no bibliography entry.  No retained line contains a figure environment, an image-inclusion command or drawing code, so no image is delivered and the package declares no asset dependency.  Editorial formatting only, in addition: an AMS \texttt{amsart} preamble that restates the article's own declarations and macros used by the retained lines, namely the environments \texttt{assumption}, \texttt{lemma} on separate counters, together with the standard packages those lines need.  The retained environments print in this document as Assumption 1, Lemma 1 in order of appearance, whereas the article prints the same items with the numbering of its own full document.  The complete native source of the article as distributed in the arXiv e-print for this exact version is bundled unchanged as \texttt{sources/203421/source0.tex}.
\begin{assumption}\label{assumption}
From now on, we make the following assumption:
\begin{enumerate}[leftmargin=2.5em]
\item The matrix $A$ is invertible.
\item  $R(iy)\leq 1$, for all real $y$.
\item $R(z)$ is analytic for $\mathrm{Re}(z)<1$.
\end{enumerate}
\end{assumption}

\begin{lemma}[Semigroup property for RKCQ]\label{lemma:Semi_group}
Consider a discrete series $\left\{ {\bf f}_k=(f_k^1,\ldots,f_k^r)^\top\right\}_{0:N}$. Then for $\alpha>0$, the RKCQ operators satisfy\footnote{\color{blue}In our applications, each stage value is vector-valued 
$f_k^i\in \R^d$, hence ${\bf f}_k\in \big(\R^d\big)^r\cong \R^{r\times d}$. All identities hold componentwise (apply the lemma to each component of $\R^d$).}
\[
\mathcal J_-^\alpha \circ \mathcal J_-^\alpha  {\bf f}_k
=
\mathcal J_-^{2\alpha} {\bf f}_k,\qquad \sigma\in \{-,+\}.
\]
\end{lemma}

\begin{proof}
For simplicity, we present the proof in the scalar case; the vector-valued case follows componentwise. Let  $\sigma=-$, the proof for $\sigma=+$ is similar. By definition,
\[
\mathcal J_-^\alpha {\bf f}_k
=
\sum_{n=0}^{k} W_n^{(\alpha)} {\bf f}_{k-n},
\]
where the matrix-valued weights $W_n^{(\alpha)} \in \mathbb{R}^{r\times r}$ are obtained by the generating series
\[
\sum_{n=0}^{\infty}
W_n^{(\alpha)} z^n
=
K^{(\alpha)}\!\left(\frac{\gamma(z)}{h}\right),
\qquad
K^{(\alpha)}(s)=s^{-\alpha}.
\]
Firstly, with $K^{(\alpha)}(s)=s^{-\alpha}$, we have
$
K^{(\alpha)}(A)K^{(\alpha)}(A)
=
K^{(2\alpha)}(A)$ for 
$A=\gamma(z)/h$. Thus
\[
\left(
\sum_{n=0}^{\infty}
W_n^{(\alpha)} z^n
\right)\left(
\sum_{m=0}^{\infty}
W_m^{(\alpha)} z^m
\right)
=
K^{(2\alpha)}\!\left(\frac{\gamma(z)}{h}\right):=\sum_{\ell=0}^{\infty}
W_\ell^{(2\alpha)} z^\ell.
\]
By the Cauchy product formula for matrix-valued power series, and uniqueness of power series coefficients, we obtain
\begin{equation}\label{proof:semigroup}
     W_\ell^{(2\alpha)}
=\sum_{n =0}^{\ell}W_{n}^{(\alpha)} W_{\ell-n }^{(\alpha)},\qquad \ell \geq 0.
\end{equation}
which is valid under Assumption~\ref{assumption}. Secondly, applying $\mathcal J_-^\alpha$ to $\mathcal J_-^\alpha {\bf f}_k$ yields
\[
(\mathcal J_-^\alpha\circ \mathcal J_-^\alpha {\bf f})_k
=
\sum_{n=0}^{k}
W_n^{(\alpha)}
(\mathcal J_-^\alpha {\bf f})_{k-n}.
\]
Substituting the definition of $J_-^\alpha {\bf f}_k$ and rearranging the sums gives
\begin{align*}
(\mathcal J_-^\alpha \circ\mathcal J_-^\alpha {\bf f})_k
=
\sum_{n=0}^{k}
W_n^{(\alpha)}
\left(
\sum_{m =0}^{k-n}
W_m ^{(\alpha)} 
{\bf f}_{k-n-m}
\right)  
=^{1}\sum_{n=0}^{k}
\sum_{\ell=n}^{k} W_{n}^{(\alpha)} W_{\ell-n}^{(\alpha)} {\bf f}_{k-\ell}.
\end{align*}
where in $=^1$ we used $\ell=n+m$. On the other hand, for a fixed $\ell$ (which can range from $0$ to  $k$) the constraint  $n\leq \ell\leq k$ with  $0\leq n\leq \ell$ shows that the last quantity is equivalent to
\begin{align*}
(\mathcal J_-^\alpha \circ\mathcal J_-^\alpha {\bf f})_k=
\sum_{\ell=0}^{k} \sum_{n=0}^{\ell}W_{n}^{(\alpha)} W_{\ell-n}^{(\alpha)} {\bf f}_{k-\ell}.
\end{align*}
With identity~\eqref{proof:semigroup}, the above shows the semigroup property, which completes the proof.
\end{proof}

\end{document}
